Token log-probabilities --- computed for free during generation by any autoregressive LLM --- are widely used as zero-cost hallucination detection signals, yet the aggregation step that converts a per-token log-probability sequence into a single confidence score has never been studied systematically. We show that this choice is not arbitrary: the optimal aggregation function is determined by hallucination type, and choosing the wrong one costs up to 12 AUROC points. Through a controlled ablation isolating aggregation function as the sole experimental variable, we establish a three-step mechanism --- hallucination type determines token probability distribution shape (peaked for recall-failure, flat for imitative-falsehood), which determines aggregation function optimality (minimum vs.\ mean) --- confirmed by distributional, rank-correlation, and threshold-level evidence across LLaMA-2-7B and Mistral-7B-v0.1. On TriviaQA, minimum log-probability outperforms mean by 6--12 AUROC points; on TruthfulQA, mean outperforms minimum by 11--12 AUROC points --- both with bootstrap 95\% confidence intervals excluding zero. Additionally, raw unnormalized log-probability sum achieves the highest AUROC on factual recall (0.895--0.896), exceeding published multi-sample Semantic Entropy at one-tenth the inference cost (under our evaluation protocol). Our findings provide the first mechanism-grounded aggregation selection criterion for zero-cost hallucination detection, filling an open gap identified in recent uncertainty quantification surveys.
